\begin{theorem}[Бейкера-Джилла-Соловея]
Для некоторого оракула $A$ выполнено $\p^A = \np^A$, а для некоторого другого оракула $B$ выполнено $\p^B \not= \np^B$.
\end{theorem}

 Идея состоит в том, что за счет использования оракула может расшириться как класс $\p$, так и класс $\np$. Теорема утверждает, что для одного оракула она заведомо расширяется одинаково, а для другого по-разному. Первая идея состоит в том, чтобы в качестве $A$ взять $\np$-полную задачу. Однако это не сработает: класс $\p$ с таким оракулом включит в себя и $\np$, и $\conp$, и еще некоторые задачи, но $\np$ тоже расширится, причем в некоторых предположениях значительно сильнее (это есть в книжке Д.В.Мусатова в разделе $6.4$).

 \begin{proof}
В качестве оракула $A$ возьмем язык $\mathsf{EXPCOM} = \{(M, x, 1^t)| M(x) = 1$ и $M(x)$ работает не дольше $2^t$ шагов $\}$. Неформально говоря, этот оракул позволяет проводить экспоненциальные вычисления за $1$ шаг. Равенство следует из того, что $\mathsf{EXP} \subset \p^{\mathsf{EXPCOM}} \subset \np^{\mathsf{EXPCOM}} \subset \mathsf{EXP}$. Первое вложение выполнено, исходя из самой сути  оракула: если $A \in \mathsf{EXP}$ и принадлежность к $A$ распознается машиной $M$ за время $2^{p(n)}$, то достаточно запросить оракул о тройке $(M, x, 1^{p(n)})$. Второе вложение выполнено для любого оракула. Наконец, третье вложение выполнено, поскольку экспоненциального времени хватит как для перебора всех ветвей алгоритма, так и для непосредственного вычисления ответов на запросы к оракулу.\\

С оракулом $B$ ситуация не такая явная. Идея состоит в том, чтобы построить $B$ постепенно, чтобы гарантировать ошибку для всех возможных полиномиальных алгоритмов. Для начала заметим, что при любом $B$ язык $U_B = \{1^n \ | \ \exists x \in B \cap \{0, 1\}^n \}$ лежит в $\np^B$. Действительно, можно недетерминированно выбрать слово $x \in \{0, 1\}^n$ и запросиь оракул о принадлежности к нему этого слова. Мы будем строить $B$ так, чтобы гарантировать $U_B \not\in \p^B$.\\

Построение языка разделим на стадии. На каждой стадии определяется принадлежность к $B$ только конечного числа слов. Более того, после каждой стадии для каждой длины либо про все слова известно, лежат ли они в $B$, либо ни про одно не известно. На стадии $k$ выберем минимальную длину $m$, такую что про слова этой длины еще ничего не известно, и запустим машину $M_k$ ($k$-ую машину Тьюринга в такой нумерации, что у каждой машины Тьюринга бесконечно много номеров) на входе $1^m$ на $\frac{2^m}{10}$ шагов. Если в процессе выполнения машина делает запросы к оракулу про слова, принадлежность которых уже установлена, отвечаем соответственно. Если она делает запросы про новые слова, отвечаем <<нет>> и помечаем, что эти слова не лежат в $B$. Наконец, если машина дает некоторый ответ, делаем так, чтобы она ошиблась: если ответ <<да>>, то помечаем, что все слова длины $m$ не лежат в $B$ (это можно сделать, т.к. перед запуском ни одно слово длины $m$ не было помечено, а в процессе могло помечаться только отсутствие в $B$), а если ответ <<нет>>, то выбираем некоторое слово длины $m$, не рассмотренное в ходе запуска (такое есть, поскольку всего слов $2^m$, а машина запросила не более $\frac{2^m}{10}$), и говорим, что оно лежит в $B$. Далее про все нерассмотренные слова длины $m$, а также всех длин, про которые машина делала запросы, говорим, что они не принадлежат $B$, и переходим к следующей стадии.\\

Ясно, что в ходе этого процесса про каждое слово рано или поздно будет сказано, лежит ли оно в $B$ или нет (именно для этого мы выбирали минимальную длину $m$; также заметим, что из конструкции получилось, что $m \geq k$). Докажем, что для полученного $B$ действительно $U_B \not\in \p^B$. Пусть это не так, и $U_B$ распознается некоторым алгоритмом за время $p(n)$. Поскольку у каждого алгоритма бесконечно много номеров, выберем такой номер $k$, что $\frac{2^m}{10} > p(m) \ \forall \ m \geq k$. В таком случае на стадии $k$ построения $B$ мы запустили $M_k$ на $1^m$ для некоторого $m \geq k$. Поскольку время работы $M_k$ меньше времени, на которое мы ее запустим, она должна выдать некоторый ответ. Но $B$ было построено таким образом, чтобы этот ответ был неправильным. Значит, никакого полиномиального алгоритма для распознавания $U_B$ нет, ЧТД.
 \end{proof}